% =====================================================================
% §1 INTRODUCTION — 0.75 pp, no figure.
% Plan per STRUCTURE.md: P1 hook-a, P2 hook-b, P3 chain+contributions,
% P4 scope+organisation.
% This file holds the POINTS TO BE EXPRESSED per paragraph, as bullets.
% Prose gets written here only after the points are agreed.
% =====================================================================
\section{Introduction}
\label{sec:intro}
Deep Neural Networks (DNNs) have become the dominant workload class across computing. They underpin the vast majority of AI-based applications, and DNN inference has become infrastructure, running continuously from edge devices to warehouse-scale data centers~\cite{jouppiInDatacenterPerformanceAnalysis2017}. The latency and energy budget of DNN accelerators is therefore a first-order design metric, and a primary axis of competition among chip designers and manufacturers. Weight movement is what drives this budget in conventional accelerators~\cite{horowitz11ComputingsEnergy2014,szeEfficientProcessingDeep2017}. In-Memory Computing (IMC) eliminates weight traffic by keeping weights in place and performing Multiply-Accumulate (MAC) operations locally. This alleviates the weight-movement bottleneck, improving performance while reducing the communication energy budget and launching a new generation of IMC DNN accelerators~\cite{sebastianMemoryDevicesApplications2020,shafieeISAACConvolutionalNeural2016,chiPRIMENovelProcessinginMemory2016}. However, no single array can hold all the weights of a full DNN, making a tiled architecture inevitable. For example, one ResNet-50 block requires 736 crossbars (CBs) and a full network requires over ten thousand. Because ADC peripherals have a large area and power overhead, they are shared across the area of a tile. This limits the number of CBs that fit in a tile, so the computation is split among hundreds or even thousands of tiles connected through a Network-on-Chip (NoC)~\cite{shafieeISAACConvolutionalNeural2016}. 


The consequence is that IMC does not eliminate communication; it relocates it. What remains mobile is activation scatter traffic and the dominant partial-sum reduction traffic: with default packing, roughly $75$--$96\%$ of a block's mobile bytes are reduction traffic that converge on a few accumulator tiles via the on-chip network (Table~\ref{tab:workloads}). The bottleneck that governs the accelerator's performance and energy becomes communication, and the design focus shifts from being computation-centric to communication-centric. Communication on IMC DNN accelerators is not set by a single design choice, nor within a single design layer, but by a chain of dependent choices taken at different layers.

This paper presents a cross-layer design space exploration spanning three main layers---packing, mapping and routing---in which the design choices available at each layer are evaluated both for their performance impact within that layer and for their impact on the layers downstream. A delay- and inference-time-aware design flow is proposed and is shown to significantly reduce both average and tail delays, the latter measured as the p99 latency that governs inference time. The contributions are:
\begin{itemize}
\item an attribution of communication delay to the layer that controls
      it: an invariant port floor (peak load on any tile
      port) fixed by the packing orientation (how a layer's rows and columns
      are grouped within a tile) and a peak link load fixed by the placement,
      with the binding term determining which layer holds control;
\item identifying the conditional couplings between adjacent design layers: what a design layer
can move or gain is set by the regime fixed by the layer above it;
\item a design flow that follows from the attribution, together with a mapping-layer design-time objective (escape slope) that lowers delay at the knee at no communication-cost overhead;
\item a bounding result: no offline quantity, over a registered search,
      predicts which placements benefit from adaptive selection; this
      establishes the limit of design-time optimization and motivates
      online adaptivity.
\end{itemize}

Results show that the packing gain dominates, up to $107\times$ in p99 at
matched load; the mapping refinement adds up to $67\%$ in p99 at zero
communication-cost overhead; and spatially adaptive selection adds up to
$2.39\times$ over a local selection heuristic. Each gain was measured with
the upstream layers already at their best, so downstream layers add to the
upstream ones rather than trading against them.